\chapter{Wielandt Bound: Supporting Results}
\label{ch:wielandt_supporting_results}

This appendix supports Chapter~\ref{ch:wielandt}. It collects exact-word,
cumulative-span, augmentation, blocking, complementary-gap, and one-step
padding results used by the main quantitative argument.

\section{Exact-word and block-injectivity support}
\label{sec:app_wielandt_exact_word_support}

This section proves the exact-word identities and low-dimensional facts used in
Lemma~\ref{lem:wordspan_blk_inj}, Theorem~\ref{thm:wordspan_blk_inj_succ},
Theorem~\ref{thm:lemma2b_assembly}, and
Theorem~\ref{thm:tp_primitive_irred_block_injective}.

\begin{lemma}[Word span products]\label{lem:word_span_mul}
    \lean{MPSTensor.wordSpan_mul_le}
    \leanok
    \uses{def:word_span}
    One has $S_m(A)S_n(A)\subseteq S_{m+n}(A)$:
    the product of word spans at lengths $m$ and $n$
    is contained in the word span at length $m+n$.
\end{lemma}

\begin{proof}\leanok
    Every product $A^{w_1}A^{w_2}$ with $|w_1|=m$
    and $|w_2|=n$ equals $A^{w_1w_2}$, where
    $|w_1w_2|=m+n$.
\end{proof}

\begin{lemma}[Zero-length word span]
    \label{lem:wordspan_zero_ne_top}
    \lean{MPSTensor.wordSpan_zero_ne_top_of_two_le}
    \leanok
    \uses{def:word_span}
    If $D\ge2$, then $S_0(A)\neq\MN{D}$.
\end{lemma}

\begin{proof}\leanok\uses{def:word_span}
    The zero-length products consist only of the identity, so
    $S_0(A)=\spn_{\C}\{\Id\}$ has dimension~$1$.
    Since $\dim\MN{D}=D^2\ge4$, this span cannot equal $\MN{D}$.
\end{proof}

\begin{corollary}[$0$-block injectivity]
    \label{cor:blk_inj_zero_dim_one}
    \lean{MPSTensor.bondDim_eq_one_of_isNBlkInjective_zero}
    \leanok
    \uses{lem:wordspan_blk_inj}
    If $D\ge1$ and $A$ is $0$-block injective, then $D=1$.
\end{corollary}

\begin{proof}\leanok\uses{lem:wordspan_blk_inj, lem:wordspan_zero_ne_top}
    By Lemma~\ref{lem:wordspan_blk_inj}, $0$-block injectivity gives
    $S_0(A)=\MN{D}$. Lemma~\ref{lem:wordspan_zero_ne_top} excludes
    $D\ge2$, so the nonzero bond dimension must be~$1$.
\end{proof}

\begin{lemma}[Two-sided span from one nonzero matrix]
    \label{lem:two_sided_nonzero_matrix_span}
    \lean{Matrix.span_range_mul_nonzero_mul_eq_top}
    \leanok
    If $C\in\MN{D}$ is nonzero, then
    \begin{align}
        \spn_{\C}\{RCS:R,S\in\MN{D}\}
         & =\MN{D}.
        \label{eq:wld_two_sided_nonzero_span}
    \end{align}
    This is the matrix-factorization input used in~%
    \cite[Lemma~3]{PerezGarcia2007Matrix}.
\end{lemma}

\begin{proof}\leanok
    Choose a nonzero entry $C_{pq}$. Then every matrix unit $E_{ij}$
    is a scalar multiple of $E_{ip}CE_{qj}$, so the span in
    (\ref{eq:wld_two_sided_nonzero_span}) contains the standard matrix basis.
\end{proof}

\begin{lemma}[Exact words around a nonzero middle matrix]
    \label{lem:block_injective_two_sided_word_span}
    \lean{MPSTensor.span_range_evalWord_mul_nonzero_mul_evalWord_eq_top}
    \leanok
    \uses{def:l_blk_injective, lem:two_sided_nonzero_matrix_span}
    If $A$ is $L$-block injective and $X\neq0$, then
    \begin{align}
        \spn_{\C}\{A^uXA^v:|u|=|v|=L\}
         & =\MN{D}.
        \notag
    \end{align}
\end{lemma}

\begin{proof}\leanok\uses{lem:two_sided_nonzero_matrix_span}
    By $L$-block injectivity, the length-$L$ word products span
    $\MN{D}$. Bilinearity of $(R,S)\mapsto RXS$ reduces the claim to
    Lemma~\ref{lem:two_sided_nonzero_matrix_span}, with $R$ and $S$ chosen
    among those word products.
\end{proof}

\section{Cumulative-span and spectral linear algebra}
\label{sec:app_wielandt_cumulative_support}

This section supplies the finite-dimensional stabilization and spectral facts
used in Theorem~\ref{thm:cumulative_span_sharp},
Theorem~\ref{thm:eigenvector_spreading}, and
Theorem~\ref{thm:wielandt_general}.

The next two lemmas connect irreducible matrix actions with the cumulative word
span used throughout Chapter~\ref{ch:wielandt}.

\begin{remark}[Burnside input for cumulative spanning]\leanok
    \uses{thm:irreducible_tensor_to_action, thm:burnside_matrix,
        lem:algspan_cumulative_span}
    The passage from irreducible tensors to full algebra spanning first passes
    through the irreducible action on $\C^D$. Burnside's theorem then says that
    the resulting irreducible subalgebra of $\MN{D}$ is all of $\MN{D}$. Once
    the algebra span is full, finite-dimensionality gives a single cumulative
    span that already equals $\MN{D}$.
\end{remark}

\begin{lemma}[Irreducible action gives an irreducible tensor]
    \label{lem:irreducible_action_to_tensor}
    \lean{MPSTensor.isIrreducibleTensor_of_isIrreducibleAction}
    \leanok
    \uses{def:irreducible_action, def:irreducible_tensor}
    If the letters of $A$ act irreducibly on $\C^D$, then $A$ has no nontrivial
    invariant orthogonal projection.
\end{lemma}

\begin{proof}
    \leanok
    If $P$ were a nontrivial invariant orthogonal projection, then its range
    would be invariant under every letter of $A$.
    Irreducibility of the action forces this range to be either $0$ or all of
    $\C^D$, so $P$ must be $0$ or $\Id$, a contradiction.
\end{proof}

\begin{lemma}[The algebra span reaches a finite cumulative span]
    \label{lem:algspan_cumulative_span}
    \lean{MPSTensor.exists_cumulativeSpan_eq_top_of_algSpan_eq_top}
    \leanok
    \uses{def:alg_span, def:cumulative_span}
    If the generated unital algebra satisfies $\algspn(A)=\MN{D}$, then there
    exists $N$ such that $T_N(A)=\MN{D}$.
\end{lemma}

\begin{proof}
    \leanok
    Every element of $\algspn(A)$ lies in some cumulative span $T_n(A)$:
    the generators lie in $T_1(A)$, scalar multiples of the identity lie in
    $T_0(A)$, and the family $\{T_n(A)\}_{n\ge0}$ is closed under addition
    and under multiplication with lengths adding.
    Since $\MN{D}$ is finite-dimensional, the ascending chain
    $T_0(A)\subseteq T_1(A)\subseteq\cdots$ stabilizes, so one can choose a
    single $N$ with $T_N(A)=\MN{D}$.
\end{proof}

\begin{definition}[Fitting decomposition]\label{def:fitting_decomp}
    \lean{Wielandt.FittingDecomposition}
    \leanok
    A \emph{Fitting decomposition} of a linear endomorphism
    $f:V\to V$ on a finite-dimensional vector space over an
    algebraically closed field consists of:
    \begin{enumerate}
        \item $f$ is nilpotent on the generalized $0$-eigenspace,
        \item $f$ is invertible on each generalized $\mu$-eigenspace
              for $\mu\neq0$,
        \item the generalized eigenspaces span~$V$,
        \item the generalized eigenspaces are linearly independent.
    \end{enumerate}
\end{definition}

\begin{theorem}[Fitting decomposition exists]\label{thm:fitting_decomp}
    \lean{Wielandt.fittingDecomposition}
    \leanok
    \uses{def:fitting_decomp}
    Every linear endomorphism on a finite-dimensional vector space
    over an algebraically closed field admits a Fitting decomposition.
\end{theorem}

\begin{proof}
    \leanok
    Nilpotency on the zero generalized eigenspace follows from
    the definition of generalized eigenspaces.
    Invertibility on nonzero generalized eigenspaces follows
    because $f-\mu$ is nilpotent there, so
    $f=\mu(1-(1-f/\mu))$ is invertible.
    Spanning and independence are standard results for
    generalized eigenspaces over algebraically closed fields.
    In~\cite{Sanz2010Quantum} and \cite[Lemma~6.3(b)]{Wolf2012Quantum},
    the Jordan normal form is used directly.
    The argument is phrased in terms of generalized eigenspaces instead.
\end{proof}

\begin{theorem}[Nilpotent power bound]\label{thm:nilpotent_pow_bound}
    \lean{Wielandt.nilpotent_pow_eq_zero_of_finrank}
    \leanok
    If $f$ is nilpotent on a space of dimension~$n$, then $f^n=0$.
\end{theorem}

\begin{proof}\leanok
    A nilpotent endomorphism on an $n$-dimensional space has
    nilpotency index at most~$n$.
\end{proof}

\begin{theorem}[Cumulative span reaches $\MN{D}$ --- explicit bound]\label{thm:cumulative_eq_top_bound}
    \lean{MPSTensor.cumulativeSpan_eq_top_of_isNormal_bound}
    \leanok
    \uses{def:cumulative_span, def:normal}
    Let $D>0$. If $A$ is normal, then
    $T_{D^2}(A)=\MN{D}$.
\end{theorem}

\begin{proof}
    \leanok
    \uses{thm:cumulative_strict_mono, lem:cumulative_finrank_le}
    By Theorem~\ref{thm:cumulative_strict_mono}, if
    $T_n\neq T_{n+1}$, then $\dim T_{n+1}>\dim T_n$.
    Since $\dim T_n\le D^2$ by Lemma~\ref{lem:cumulative_finrank_le},
    after at most $D^2$ steps either $T_n=\MN{D}$ or $T_n$ has stabilised.
    Stabilisation before reaching $\MN{D}$ contradicts normality, which
    requires $S_{L_0}=\MN{D}$ for some $L_0$ and hence
    $T_{L_0}=\MN{D}$.
\end{proof}

\begin{theorem}[Nonzero trace product]\label{thm:nonzero_trace_word}
    \lean{MPSTensor.exists_nonzero_trace_word}
    \leanok
    \uses{def:normal}
    Let $D>0$. If $A$ is normal, then there exists a word $w$ of length at
    most $D^2$ such that $\tr(A^w)\neq0$.
    This is the weaker $D^2$ corollary of
    \cite[Lemma~1]{Sanz2010Quantum}. The sharp source statement is
    Theorem~\ref{thm:nonzero_trace_word_primitive_sharp}.
\end{theorem}

\begin{proof}
    \leanok
    \uses{thm:cumulative_eq_top}
    By Theorem~\ref{thm:cumulative_eq_top}, $\Id\in T_{D^2}(A)$, so $\Id$
    is a linear combination of word products of length at most~$D^2$.
    At least one has nonzero trace because $\tr(\Id)=D\neq0$.
\end{proof}

\begin{theorem}[Nonzero trace implies eigenvector]\label{thm:eigenvector_from_trace}
    \lean{exists_eigenvector_of_trace_ne_zero}
    \leanok
    If $M\in\MN{D}$ has $\tr(M)\neq0$, then $M$ has a nonzero eigenvalue
    $\mu\neq0$ and a corresponding eigenvector $\varphi\neq0$ satisfying
    $M\varphi=\mu\varphi$.
\end{theorem}

\begin{proof}
    \leanok
    The trace is the sum of the eigenvalues of $M$, counted with
    multiplicity, so a nonzero trace gives a nonzero eigenvalue~$\mu$.
    Choosing a nonzero vector $\varphi\in\ker(M-\mu\Id)$ gives
    $M\varphi=\mu\varphi$.
\end{proof}

\begin{theorem}[Eigenvalue extraction]\label{thm:eigenvalue_extraction}
    \lean{MPSTensor.exists_word_eigenvector}
    \leanok
    \uses{def:normal}
    Let $D>0$. If $A$ is normal, then there exist a word $w_0$ with
    $|w_0|\le D^2$, a nonzero scalar $\mu\neq0$, and a nonzero vector
    $\varphi\neq0$ such that $A^{w_0}\varphi=\mu\varphi$.
\end{theorem}

\begin{proof}
    \leanok
    \uses{thm:nonzero_trace_word, thm:eigenvector_from_trace}
    By Theorem~\ref{thm:nonzero_trace_word}, there is a word $w_0$ with
    $\tr(A^{w_0})\neq0$ and $|w_0|\le D^2$.
    By Theorem~\ref{thm:eigenvector_from_trace}, $A^{w_0}$ has a nonzero
    eigenvector.
\end{proof}

\begin{theorem}[Vector span stabilisation]\label{thm:vector_span_stable}
    \lean{MPSTensor.cumulativeVectorSpan_stable}
    \leanok
    \uses{def:cumulative_vector_span}
    If $K_n(A,\varphi)=K_{n+1}(A,\varphi)$, then
    $K_m(A,\varphi)=K_n(A,\varphi)$ for every $m\ge n$.
\end{theorem}

\begin{proof}\leanok\uses{thm:cumulative_stable}
    The argument is the same as for Theorem~\ref{thm:cumulative_stable}:
    left multiplication by $A^i$ cannot escape the stabilised span.
\end{proof}

\begin{lemma}[Vector span from matrix span]\label{lem:vector_span_from_matrix}
    \lean{MPSTensor.cumulativeVectorSpan_eq_top_of_cumulativeSpan_eq_top}
    \leanok
    \uses{def:cumulative_vector_span, def:cumulative_span}
    If $T_N(A)=\MN{D}$ and $\varphi\neq0$, then
    $K_N(A,\varphi)=\C^D$.
\end{lemma}

\begin{proof}
    \leanok
    Every matrix $M\in\MN{D}$ lies in $T_N(A)$, so in particular every
    rank-one matrix sending $\varphi$ to any target vector~$v$ lies in
    $T_N(A)$. Applying such matrices to $\varphi$ recovers all of~$\C^D$.
\end{proof}

\section{One-step augmentation}
\label{sec:app_wielandt_augmentation}

This section proves the one-step augmentation results used in
Theorem~\ref{thm:wielandt_case2_one_step_span} and
Theorem~\ref{thm:wielandt_case3_one_step_span}.

\begin{definition}[One-step augmentation]\label{def:one_step_augment}
    \lean{MPSTensor.oneStepAugment}
    \leanok
    \uses{def:mps_tensor}
    Let $A$ be an MPS tensor with Kraus family
    $\{A^i\}_{i=0}^{d-1}$, and let $X\in\MN{D}$.
    The \emph{one-step augmentation} of $A$ by $X$ is the tensor
    $\widetilde A$ with physical dimension $d+1$ whose first Kraus operator
    is $X$ and whose remaining Kraus operators are those of $A$:
    \begin{align}
        \widetilde A^0
         & =X,
        \notag   \\
        \widetilde A^{i+1}
         & =A^i.
        \notag
    \end{align}
\end{definition}

\begin{lemma}[Kraus operators lie in the one-step span]\label{lem:kraus_mem_word_span_one}
    \lean{MPSTensor.apply_mem_wordSpan_one}
    \leanok
    \uses{def:word_span}
    Every Kraus operator $A^i$ belongs to the one-step span $S_1(A)$.
\end{lemma}

\begin{proof}
    \leanok
    This is the length-one word with letter~$i$.
\end{proof}

\begin{lemma}[Redundant one-step augmentation]\label{lem:one_step_augment_word_span}
    \lean{MPSTensor.oneStepAugment_apply_mem_wordSpan_one,
        MPSTensor.wordSpan_oneStepAugment_one,
        MPSTensor.wordSpan_oneStepAugment_eq}
    \leanok
    \uses{def:one_step_augment, def:word_span, lem:kraus_mem_word_span_one}
    Let $X\in S_1(A)$, and let $\widetilde A$ be the one-step
    augmentation of $A$ by $X$. Then
    $S_n(\widetilde A)=S_n(A)$ for every $n\ge0$.
\end{lemma}

\begin{proof}
    \leanok
    \uses{lem:kraus_mem_word_span_one}
    The new generator $X$ is already in $S_1(A)$, and each original
    generator remains in the augmented family. Hence the one-step spans
    agree. Multiplying word spans and arguing by induction on the word
    length gives equality at every length.
\end{proof}

\begin{lemma}[Normality and Kraus rank under one-step augmentation]%
    \label{lem:one_step_augment_normal_kraus_rank}
    \lean{MPSTensor.isNormal_oneStepAugment_of_mem_wordSpan_one,
        MPSTensor.krausRank_oneStepAugment}
    \leanok
    \uses{def:one_step_augment, def:normal, lem:one_step_augment_word_span}
    Let $X\in S_1(A)$, and let $\widetilde A$ be the one-step augmentation
    of $A$ by $X$. If $A$ is normal, then $\widetilde A$ is normal, and
    $\krausrk(\widetilde A)=\krausrk(A)$.
\end{lemma}

\begin{proof}
    \leanok
    \uses{lem:one_step_augment_word_span}
    Since all exact word spans are unchanged by
    Lemma~\ref{lem:one_step_augment_word_span}, eventual full word span
    for $A$ gives eventual full word span for $\widetilde A$.
    The equality of Kraus ranks is the equality of the dimensions of the
    common one-step span.
\end{proof}

\begin{theorem}[Wielandt case (2), invertible generator]%
    \label{thm:wielandt_case2_normal_generator}
    \lean{MPSTensor.wordSpan_eq_top_of_isNormal_of_isUnit,
        MPSTensor.iIndex_le_of_isNormal_of_isUnit}
    \leanok
    \uses{def:word_span, def:normal}
    Let $A$ be a normal MPS tensor with bond dimension $D>0$, and suppose
    that some Kraus operator $A^{i_0}$ is invertible. Then
    \begin{align}
        S_{D^2-\krausrk(A)+1}(A)
         & =\MN{D},
        \notag                    \\
        \iota(A)
         & \le D^2-\krausrk(A)+1.
        \notag
    \end{align}
\end{theorem}

\begin{proof}
    \leanok
    The dimensions of the exact word spans cannot decrease, because
    left multiplication by the invertible Kraus operator is injective.
    If this dimension growth stopped before the full matrix algebra was
    reached, the same argument at every later length would contradict
    normality. Since $\dim S_1(A)=\krausrk(A)$ and
    $\dim\MN{D}=D^2$, the full algebra is reached by length
    $D^2-\krausrk(A)+1$. The index bound follows from the definition
    of $\iota(A)$.
\end{proof}

\begin{corollary}[Wielandt case (2), invertible Kraus operator]%
    \label{cor:wielandt_case2_generator}
    \lean{MPSTensor.wordSpan_eq_top_of_isPrimitivePaper_of_isUnit,
        MPSTensor.iIndex_le_of_isPrimitivePaper_of_isUnit}
    \leanok
    \uses{lem:kraus_mem_word_span_one, thm:wielandt_case2_one_step_span,
        cor:wielandt_case2_one_step_index}
    Let $A$ be an MPS tensor with bond dimension $D\ge1$, satisfying
    $\sum_i(A^i)^\dagger A^i=\Id$, and suppose that $A$ is primitive in the
    sense of Definition~\ref{def:paper_primitive}.
    If some Kraus operator $A^{i_0}$ is invertible, then
    \begin{align}
        S_{D^2-\krausrk(A)+1}(A)
         & =\MN{D},
        \notag                    \\
        \iota(A)
         & \le D^2-\krausrk(A)+1.
        \notag
    \end{align}
\end{corollary}

\begin{proof}
    \leanok
    \uses{lem:kraus_mem_word_span_one, thm:wielandt_case2_one_step_span,
        cor:wielandt_case2_one_step_index}
    Apply the one-step case to $X=A^{i_0}$, which belongs to $S_1(A)$ by
    Lemma~\ref{lem:kraus_mem_word_span_one}.
\end{proof}

\begin{corollary}[Wielandt case (3), non-invertible Kraus operator]%
    \label{cor:wielandt_case3_generator}
    \lean{MPSTensor.wordSpan_eq_top_of_isPrimitivePaper_of_noninvertible_eigenvector,
        MPSTensor.iIndex_le_sq_of_noninvertible_eigenvector}
    \leanok
    \uses{lem:kraus_mem_word_span_one, thm:wielandt_case3_one_step_span,
        cor:wielandt_case3_one_step_index}
    Let $A$ be an MPS tensor with bond dimension $D\ge1$, satisfying
    $\sum_i(A^i)^\dagger A^i=\Id$, and suppose that $A$ is primitive in the
    sense of Definition~\ref{def:paper_primitive}.
    If some Kraus operator $A^{i_0}$ is non-invertible and has an
    eigenvector $\varphi\neq0$ with nonzero eigenvalue~$\mu$, then
    \begin{align}
        S_{D^2}(A)
         & =\MN{D},
        \notag      \\
        \iota(A)
         & \le D^2.
        \notag
    \end{align}
\end{corollary}

\begin{proof}
    \leanok
    \uses{lem:kraus_mem_word_span_one, thm:wielandt_case3_one_step_span,
        cor:wielandt_case3_one_step_index}
    Apply the one-step case to $X=A^{i_0}$, which belongs to $S_1(A)$ by
    Lemma~\ref{lem:kraus_mem_word_span_one}.
\end{proof}

\section{Blocking and fixed-length spanning}
\label{sec:app_wielandt_blocking_support}
\label{sec:rectangular_span}

This section proves the blocking and fixed-length constructions used in
Theorem~\ref{thm:wielandt_general},
Theorem~\ref{thm:bounded_injective_blocking_normal_left_canonical}, and
Theorem~\ref{thm:tp_primitive_irred_block_injective}.

\begin{theorem}[Blocking preserves normality]\label{thm:isNormal_blockTensor}
    \lean{MPSTensor.isNormal_blockTensor_of_isNormal}
    \leanok
    \uses{def:normal, def:blocked_tensor}
    If $A$ is normal and $L>0$, then the blocked tensor
    $A^{[L]}$ is also normal.
\end{theorem}

\begin{proof}\leanok
    \uses{lem:word_span_le_blockTensor, lem:word_span_mul}
    Choose $N$ with $S_N(A)=\MN{D}$. Then $\Id\in S_N(A)$.
    For any $M\in S_N(A)$, write
    $\Id=\sum_r c_rA^{w_r}$ with each $|w_r|=N$. Then
    \begin{align}
        M
         & =M\Id
        =\sum_r c_rMA^{w_r}
        \in S_N(A)S_N(A)
        \subseteq S_{2N}(A).
        \label{eq:wld_blocknormal_self_inclusion}
    \end{align}
    By Lemma~\ref{lem:word_span_mul}, iterating
    (\ref{eq:wld_blocknormal_self_inclusion}) gives
    $S_N(A)\subseteq S_{kN}(A)$ for every $k\ge1$.
    Taking $k=L$ yields
    $\MN{D}=S_N(A)\subseteq S_{NL}(A)$.
    By Lemma~\ref{lem:word_span_le_blockTensor},
    $S_{NL}(A)\subseteq S_N(A^{[L]})$, so
    $S_N(A^{[L]})=\MN{D}$.
\end{proof}

\begin{lemma}[Chunking: word span embeds into blocked word span]%
    \label{lem:word_span_le_blockTensor}
    \lean{MPSTensor.wordSpan_le_wordSpan_blockTensor}
    \leanok
    \uses{def:word_span, def:blocked_tensor}
    One has $S_{nL}(A)\subseteq S_n(A^{[L]})$ for all $n,L$.
\end{lemma}

\begin{proof}\leanok
    \uses{lem:word_span_mul}
    Every length-$nL$ word in the original alphabet
    can be split into $n$ consecutive chunks of length~$L$,
    each of which is a single blocked Kraus operator by
    (\ref{eq:mps_blocked_tensor}).
    The claim follows by induction on~$n$, using
    Lemma~\ref{lem:word_span_mul}.
\end{proof}

\begin{theorem}[Word eigenvector gives blocked eigenvector]%
    \label{thm:blockTensor_single_eigenvector}
    \lean{MPSTensor.blockTensor_single_eigenvector}
    \leanok
    \uses{def:blocked_tensor}
    If $A^{w_0}\varphi=\mu\varphi$ for a word $w_0$
    of length~$L$, then $(A^{[L]})^{j_0}\varphi=\mu\varphi$,
    where $j_0$ is the encoded blocked index corresponding
    to $w_0$.
\end{theorem}

\begin{proof}\leanok
    By (\ref{eq:mps_blocked_tensor}),
    $(A^{[L]})^{j_0}=A^{w_0}$.
\end{proof}

\begin{theorem}[Blocked fixed-length matrix spanning]\label{thm:wielandt_blocked_assembly}
    \lean{MPSTensor.wielandt_blocked_assembly}
    \leanok
    \uses{def:word_span, def:normal, def:blocked_tensor}
    Suppose $A$ has bond dimension $D>0$, is normal, and $L>0$, and we have:
    \begin{enumerate}
        \item a word $w_0$ of length $L$ with eigenvector
              $A^{w_0}\varphi=\mu\varphi$
              ($\mu\neq0$, $\varphi\neq0$),
        \item a word $\tau_0$ of length $L$ with transpose
              eigenvector
              $(A^{\tau_0})^T\psi=\nu\psi$
              ($\nu\neq0$, $\psi\neq0$),
        \item a rank-one element
              $\varphi\psi^T\in S_m(A^{[L]})$
              for some $m$.
    \end{enumerate}
    Then there exists $N$ such that $S_N(A)=\MN{D}$.

    \noindent\emph{Remark.}
    The proof gives the explicit bound
    $N=((D{-}1)+m+(D{-}1))L$; the theorem states only the
    existential conclusion.
\end{theorem}

\begin{proof}
    \leanok
    \uses{thm:isNormal_blockTensor, thm:blockTensor_single_eigenvector,
        thm:eigenvector_spreading, lem:blocking_transfer,
        thm:lemma2b_assembly}
    Set $B=A^{[L]}$, which is normal by
    Theorem~\ref{thm:isNormal_blockTensor}. The word eigenvectors transfer to
    single-index eigenvectors of $B$ and $B^T$ by
    Theorem~\ref{thm:blockTensor_single_eigenvector}. Eigenvector spreading
    and padding then give
    $H_{D-1}(B,\varphi)=H_{D-1}(B^T,\psi)=\C^D$.

    For each standard basis vector $e_j$, choose
    $R_j\in S_{D-1}(B^T)$ with $R_j\psi=e_j$. Reversing each word after
    transposition shows that $R_j^T\in S_{D-1}(B)$, and therefore
    $(\varphi\psi^T)R_j^T=\varphi e_j^T$ belongs to
    $S_{m+D-1}(B)$. Theorem~\ref{thm:lemma2b_assembly}, applied with
    $n=D-1$ and $m'=m+D-1$, now gives
    $S_{(D-1)+m+(D-1)}(B)=\MN{D}$. Finally,
    Lemma~\ref{lem:blocking_transfer} transfers this identity to
    $S_{((D-1)+m+(D-1))L}(A)=\MN{D}$.
\end{proof}

\begin{theorem}[Rank-one extraction for blocked tensor]%
    \label{thm:rankone_extraction_blockTensor}
    \lean{MPSTensor.exists_rankOne_mem_wordSpan_blockTensor}
    \leanok
    \uses{def:word_span, def:normal, def:blocked_tensor}
    Suppose $A$ has bond dimension $D>0$, is normal, and $N_0\ge1$ with
    $S_{N_0}(A)=\MN{D}$.
    Then there exist words $\sigma_0,\tau_0$ of length~$N_0$,
    nonzero vectors $\varphi,\psi$, nonzero scalars $\mu,\nu$,
    and $m\in\N$ such that:
    \begin{enumerate}
        \item $A^{\sigma_0}\varphi=\mu\varphi$,
        \item $(A^{\tau_0})^T\psi=\nu\psi$,
        \item $\varphi\psi^T\in S_m(A^{[N_0]})$.
    \end{enumerate}
\end{theorem}

\begin{proof}
    \leanok
    \uses{def:blocked_tensor, thm:isNormal_blockTensor,
        thm:eigenvector_from_trace}
    Set $B=A^{[N_0]}$.
    By (\ref{eq:mps_blocked_tensor}) and
    $S_{N_0}(A)=\MN{D}$, one has $S_1(B)=\MN{D}$.
    Hence $\Id\in S_1(B)$, so not all blocked generators can have trace
    zero.
    Choose $j_0$ with $\tr(B^{j_0})\neq0$.
    By Theorem~\ref{thm:eigenvector_from_trace},
    $B^{j_0}$ has a nonzero eigenvector $\varphi$
    with eigenvalue $\mu\neq0$.
    Writing $j_0$ as the blocked index corresponding to a word
    $\sigma_0$ of length~$N_0$ gives
    $A^{\sigma_0}\varphi=\mu\varphi$.
    Applying the same argument to the transpose blocked tensor gives
    a word $\tau_0$ of length~$N_0$, a nonzero vector $\psi$,
    and $\nu\neq0$ with $(A^{\tau_0})^T\psi=\nu\psi$.
    Since $A$ is normal and $N_0\ge1$,
    the blocked tensor $B$ is normal by
    Theorem~\ref{thm:isNormal_blockTensor}.
    Therefore there exists $M$ with $S_M(B)=\MN{D}$.
    In particular, $\varphi\psi^T\in S_M(B)$.
    Taking $m=M$ proves~(3).
\end{proof}

\begin{theorem}[Fixed-length word-span saturation via blocking]\label{thm:wielandt_lemma2b}
    \lean{MPSTensor.wielandt_lemma2b}
    \leanok
    \uses{def:word_span, def:normal}
    If $A$ is a normal MPS tensor with bond dimension $D>0$,
    then there exists $N$ such that $S_N(A)=\MN{D}$.
    This is a blocked fixed-length construction under the stated normality
    hypothesis. The quantitative result of \cite[Lemma~2(b)]{Sanz2010Quantum}
    is Theorem~\ref{thm:rank_one_noninvertible_eigenvector_primitive}.
\end{theorem}

\begin{proof}
    \leanok
    \uses{thm:rankone_extraction_blockTensor,
        thm:wielandt_blocked_assembly}
    By normality, some $N_0>0$ has $S_{N_0}(A)=\MN{D}$.
    Theorem~\ref{thm:rankone_extraction_blockTensor} produces eigenvectors
    and a rank-one element
    $\varphi\psi^T\in S_m(A^{[N_0]})$.
    Theorem~\ref{thm:wielandt_blocked_assembly} then gives
    $S_{((D-1)+m+(D-1))N_0}(A)=\MN{D}$.
\end{proof}

\section{Complementary-gap consequences}
\label{sec:app_wielandt_complementary_gap}

This section proves the complementary-gap consequences used in
Lemma~\ref{lem:primitive_implies_irreducible},
Lemma~\ref{lem:strongly_irred_from_primitive}, and
Theorem~\ref{thm:normal_from_primitive_posdef}.

\begin{theorem}[Primitive overlap convergence]\label{thm:prim_overlap_one}
    \lean{MPSTensor.IsPrimitiveMPS.overlap_tendsto_one}
    \leanok
    \uses{def:primitive_mps, def:mpv_overlap}
    If $A$ is a primitive MPS tensor, then $O_{AA}(N)\to1$.
\end{theorem}

\begin{proof}\leanok\uses{thm:primitive_overlap}
    The claim follows from the self-overlap convergence under a
    complementary transfer-map gap
    (Theorem~\ref{thm:primitive_overlap}).
\end{proof}

\begin{theorem}[Peripheral primitivity and irreducibility give normalized overlap]
    \label{thm:peripheral_primitive_irreducible_overlap_one}
    \lean{MPSTensor.%
        overlap_tendsto_one_of_peripheralPrimitive_of_irreducible}
    \leanok
    \uses{def:primitive_channel, def:irreducible_tensor, def:mpv_overlap}
    Let $A$ have positive bond dimension and satisfy
    $\sum_i(A^i)^\dagger A^i=\Id$.
    If $A$ is tensor-irreducible and the transfer map is primitive, then
    $O_{AA}(N)\to1$.
\end{theorem}

\begin{proof}\leanok\uses{thm:prim_overlap_one}
    Irreducible Perron--Frobenius theory supplies a nonzero positive fixed
    point. Peripheral primitivity gives a strict spectral-radius bound on the
    complementary transfer map. Thus $A$ is a primitive MPS tensor, and
    Theorem~\ref{thm:prim_overlap_one} applies.
\end{proof}

\begin{lemma}[Primitive MPS tensors have nonzero word products]\label{lem:prim_nonzero_word}
    \lean{MPSTensor.exists_nonzero_evalWord_of_isPrimitiveMPS}
    \leanok
    \uses{def:primitive_mps, def:eval_word}
    If $A$ is a primitive MPS tensor,
    then for every $n\ge0$ there exists a word $w$ of length
    \emph{exactly}~$n$ with $A^w\neq0$.
\end{lemma}

\begin{proof}\leanok
    The PSD fixed point $\rho$ of $\E_A$ satisfies
    $\E_A^n(\rho)=\rho\neq0$.
    Since
    $\E_A^n(\rho)=\sum_{|w|=n}A^w\rho(A^w)^\dagger$,
    at least one summand is nonzero, hence some $A^w\neq0$.
\end{proof}

\begin{lemma}[Iterated transfer does not vanish]\label{lem:transfer_pow_ne_zero}
    \lean{MPSTensor.transferMap_pow_ne_zero_of_isPrimitiveMPS}
    \leanok
    \uses{def:primitive_mps, def:transfer_map}
    If $A$ is a primitive MPS tensor, then
    $\E_A^n\neq0$ for all $n\ge0$.
\end{lemma}

\begin{proof}\leanok\uses{lem:prim_nonzero_word}
    By Lemma~\ref{lem:prim_nonzero_word}, there exists a word~$w$ of
    length~$n$ with $A^w\neq0$.
    The expansion
    $\E_A^n(\Id)=\sum_{|u|=n}A^u(A^u)^\dagger$
    contains the nonzero positive semidefinite summand
    $A^w(A^w)^\dagger$, so $\E_A^n(\Id)\neq0$.
\end{proof}

\begin{lemma}[Fixed-point uniqueness from a complementary transfer-map gap]%
    \label{lem:fixedpoint_unique_spectral}
    \lean{MPSTensor.IsPrimitiveMPS.fixedPoint_unique}
    \leanok
    \uses{def:primitive_mps}
    If $A$ is a primitive MPS tensor
    with PSD fixed point $\rho$,
    then every fixed point $\sigma$ of $\E_A$ is proportional
    to $\rho$: $\sigma=\frac{\tr(\sigma)}{\tr(\rho)}\rho$.
\end{lemma}

\begin{proof}\leanok
    Set $\sigma'=\sigma-\frac{\tr(\sigma)}{\tr(\rho)}\rho$.
    Then $\tr(\sigma')=0$ and $(\E_A-P_\rho)(\sigma')=\sigma'$.
    If $\sigma'\neq0$, then $1$ is an eigenvalue of $\E_A-P_\rho$,
    contradicting the complementary transfer-map gap
    $\rho_{\spec}(\E_A-P_\rho)<1$.
\end{proof}

\begin{lemma}[Complement powers tend to zero]%
    \label{lem:complement_pow_zero}
    \lean{MPSTensor.IsPrimitiveMPS.complement_pow_tendsto_zero}
    \leanok
    \uses{def:primitive_mps}
    If $A$ is a primitive MPS tensor with PSD fixed point $\rho$,
    then $(\E_A-P_\rho)^n\to0$ in operator norm.
\end{lemma}

\begin{proof}\leanok
    The spectral radius of $\E_A-P_\rho$ is strictly less
    than $1$ by the complementary transfer-map gap hypothesis.
    Apply the general result that powers of an operator
    with spectral radius $<1$ tend to zero.
\end{proof}

\begin{lemma}[Primitive tensors with positive-definite fixed point have irreducible transfer map]%
    \label{lem:irreducible_map_from_primitive}
    \lean{MPSTensor.isIrreducibleMap_of_isPrimitiveMPS_of_posDef}
    \leanok
    \uses{def:primitive_mps}
    If $A$ is a primitive MPS tensor with positive-definite fixed
    point~$\rho$, then the transfer map $\E_A$ is irreducible.

    The proof uses fixed-point uniqueness
    (Lemma~\ref{lem:fixedpoint_unique_spectral}):
    every PSD fixed point is proportional to $\rho$,
    so Wolf's criterion for irreducibility
    (a positive-definite fixed point together with uniqueness of positive
    fixed points implies irreducibility) applies directly.
\end{lemma}

\begin{proof}
    \leanok
    \uses{lem:fixedpoint_unique_spectral}
    By Lemma~\ref{lem:fixedpoint_unique_spectral},
    if $\sigma\ge0$ and $\E_A(\sigma)=\sigma$, then
    $\sigma=\frac{\tr(\sigma)}{\tr(\rho)}\rho$.
    Since $\rho>0$, Wolf's uniqueness criterion for
    irreducibility applies.
\end{proof}

\section{Identity in the one-step span}
\label{sec:app_wielandt_identity_in_span}

This section proves the one-step padding conversions used in
Theorem~\ref{thm:algspan_to_normal}. It also records why cumulative spanning
alone does not imply normality.

\begin{remark}[Identity-in-the-one-step-span condition]
    \label{rem:aperiodicity}\leanok
    The \emph{one-step padding condition} is $\Id\in S_1(A)$: the identity
    matrix lies in the span of the Kraus operators
    $\{A^i\}_{i=0}^{d-1}$. This is an additional hypothesis. The
    normalization $\sum_i(A^i)^\dagger A^i=\Id$ gives a quadratic identity
    and does not imply $\Id\in S_1(A)$.
    Theorem~\ref{thm:normal_from_primitive_posdef} obtains normality from
    a positive-definite fixed point and a complementary gap without this
    padding condition.
\end{remark}

\begin{remark}[Counterexample: cumulative spanning does not imply normality]%
    \label{rem:cumspan_counterexample}\leanok
    In $\MN{2}$, let $A^0=E_{12}$ and $A^1=E_{21}$ be the off-diagonal
    matrix units.
    Then $\algspn(A)=\MN{2}$ and $T_2(A)=\MN{2}$, but the word spans
    alternate: $S_n(A)$ contains only off-diagonal matrices for odd~$n$
    and only diagonal matrices for even~$n$.
    Hence $S_n(A)\neq\MN{2}$ for every~$n$, and $A$ is not normal.
    The one-step padding condition fails because
    $\Id\notin S_1(A)=\spn_{\C}\{E_{12},E_{21}\}$.
\end{remark}

\begin{theorem}[Word-span monotonicity from one-step padding]%
    \label{thm:wordspan_mono_aperiodic}
    \lean{MPSTensor.wordSpan_mono_succ_of_one_mem_wordSpan_one}
    \leanok
    \uses{def:word_span}
    If $\Id\in S_1(A)$, then $S_n(A)\le S_{n+1}(A)$ for every~$n$.
\end{theorem}

\begin{proof}\leanok\uses{lem:word_span_mul}
    For any $M\in S_n(A)$, one has
    $M=M\Id\in S_n(A)S_1(A)\subseteq S_{n+1}(A)$ by
    Lemma~\ref{lem:word_span_mul}.
\end{proof}

\begin{theorem}[Cumulative span equals word span under one-step padding]%
    \label{thm:cumspan_eq_wordspan}
    \lean{MPSTensor.cumulativeSpan_eq_wordSpan_of_one_mem_wordSpan_one}
    \leanok
    \uses{def:cumulative_span, def:word_span}
    If $\Id\in S_1(A)$, then $T_n(A)=S_n(A)$ for every~$n$.
\end{theorem}

\begin{proof}\leanok\uses{thm:wordspan_mono_aperiodic}
    By Theorem~\ref{thm:wordspan_mono_aperiodic}, the word spans form the
    chain $S_0(A)\le S_1(A)\le\cdots\le S_n(A)$.
    Therefore their sum $T_n(A)=S_0(A)+\cdots+S_n(A)$ equals the largest
    term $S_n(A)$.
\end{proof}

\begin{theorem}[Cumulative spanning plus one-step padding implies normality]%
    \label{thm:cumspan_to_normal}
    \lean{MPSTensor.isNormal_of_cumulativeSpan_eq_top_of_aperiodic}
    \leanok
    \uses{def:normal, def:cumulative_span, def:word_span}
    If $T_N(A)=\MN{D}$ and $\Id\in S_1(A)$, then $A$ is normal.
\end{theorem}

\begin{proof}\leanok\uses{thm:cumspan_eq_wordspan}
    Theorem~\ref{thm:cumspan_eq_wordspan} gives
    $S_N(A)=T_N(A)=\MN{D}$.
\end{proof}
